\documentclass{article}
\usepackage[utf8]{inputenc}
\usepackage{amsmath,amssymb}
\usepackage[most]{tcolorbox}
\usepackage{geometry}
\geometry{margin=2.5cm}

\newcommand{\step}[1]{\par\noindent\hangindent=3.5em\hangafter=1%
  \makebox[3.5em][l]{\textbf{#1.}}}
\newcommand{\steptag}[1]{\hfill{\small\textsf{[#1]}}}

\newtcolorbox{proofbox}[1]{
  colback=gray!5, colframe=gray!60,
  fonttitle=\bfseries, title={#1},
  breakable, enhanced,
  left=4pt, right=4pt, top=4pt, bottom=4pt,
  before skip=10pt, after skip=10pt
}

\begin{document}

\begin{proofbox}{There are universal C\_out < infinity and e\_out > 0 such t...}
\step{1} There are universal C\_out < infinity and e\_out > 0 such that, whenever R,P are t-projections in a finite-dimensional extended t-C*-algebra, R is nonvanishing, both subordination errors of P to R are at most t, v:B->S\textasciicircum{}A\_P is an extended t-isomorphism, A\_R = S\textasciicircum{}A\_R, P\textasciicircum{}R = Co\textasciicircum{}A\_R(P), and t <= e\_out, the explicitly defined map T = Co\textasciicircum{}\{A\_R\}\_\{P\textasciicircum{}R\} o Co\textasciicircum{}A\_R o v : B->S\textasciicircum{}\{A\_R\}\_\{P\textasciicircum{}R\} is an extended C\_out*t-isomorphism and T\_n = I\_n tensor T for every n. \steptag{validated} \\[6pt]
\step{1.1} Nested-corner geometry and target algebra ledger. There are universal constants C\_nest,C\_ca and a positive universal threshold such that, with a=P\textasciicircum{}R=Co\textasciicircum{}A\_R(P), the maps W\_n=Co\textasciicircum{}\{A\_R\}\_\{I\_n tensor a\} composed with Co\textasciicircum{}A\_\{I\_n tensor R\} carry M\_n tensor S\textasciicircum{}A\_P into M\_n tensor S\textasciicircum{}\{A\_R\}\_a, satisfy ||W\_n(X)-X||<=C\_nest*t*||X||, and dim S\textasciicircum{}A\_P=dim S\textasciicircum{}\{A\_R\}\_a. Moreover S\textasciicircum{}A\_P, A\_R, and S\textasciicircum{}\{A\_R\}\_a have their compressed extended approximate C*-algebra structures with defects bounded by universal multiples of t. \steptag{validated} \\[6pt]
\step{1.1.1} Specialize the validated externals lem-maincb-nested-corner-comparison and lem-maincb-nested-corner-dimension-transport to Q=P. The two left/right subordination assumptions then supply all four repeated hypotheses. Thus a=Co\textasciicircum{}A\_R(P) is a C\_nest*t-projection in A\_R; for every n the forward comparison W\_n=Co\textasciicircum{}\{A\_R\}\_\{I\_n tensor a\} composed with Co\textasciicircum{}A\_\{I\_n tensor R\} maps M\_n tensor S\textasciicircum{}A\_P to M\_n tensor S\textasciicircum{}\{A\_R\}\_a and satisfies ||W\_n(X)-X||<=C\_nest*t*||X||; and dim S\textasciicircum{}A\_P=dim S\textasciicircum{}\{A\_R\}\_a. These conclusions use the two externals by their exact registered names and no re-proof of them. \steptag{validated} \\[4pt]
\step{1.1.2} Nonvanishing and corner-algebra check. A t-projection in an extended t-C*-algebra has norm at most 2 for a universal small threshold, since (1-t)||H||\textasciicircum{}2<=||H\textasciicircum{}2||<=||H||+t. If P were in the vanishing alternative of def-delta-projection, then ||P||=O(t), and the two operator comparisons in def-compressed-corner would make the idempotent Co\textasciicircum{}A\_P have operator norm O(t)<1; lem-compcb-amplified-compression-identities then forces Co\textasciicircum{}A\_P=0. This contradicts S\textasciicircum{}A\_P!=0, because the bijection v has nonzero unital source B. Hence P is nonvanishing. Applying lem-compcb-corner-algebra to P and to the given nonvanishing R makes S\textasciicircum{}A\_P and A\_R=S\textasciicircum{}A\_R extended 2*C\_ca*t-C*-algebras. The dimension equality in node 1.1.1 gives S\textasciicircum{}\{A\_R\}\_a!=0. If the C\_nest*t-projection a were vanishing in A\_R, the identical compression-operator argument, now with total error (C\_nest+2*C\_ca)t and using lem-compcb-amplified-compression-identities inside A\_R, would force Co\textasciicircum{}\{A\_R\}\_a=0, a contradiction. Thus a is nonvanishing, and a second application of lem-compcb-corner-algebra makes S\textasciicircum{}\{A\_R\}\_a an extended C\_ca*(C\_nest+2*C\_ca)*t-C*-algebra. All thresholds are positive universal finite minima. \steptag{validated} \\[4pt]
\step{1.2} Exact amplified formula, range, and involution. Put R\_n=I\_n tensor R and a\_n=I\_n tensor a. Tensor functoriality for linear maps and lem-compcb-amplified-compression applied first in A and then in A\_R give I\_n tensor T=Co\textasciicircum{}\{A\_R\}\_\{a\_n\} composed with Co\textasciicircum{}A\_\{R\_n\} composed with v\_n=W\_n composed with v\_n. The same external identifies M\_n tensor S\textasciicircum{}A\_P and M\_n tensor S\textasciicircum{}\{A\_R\}\_a with the corresponding amplified compression ranges, so the displayed map is linear and has the asserted target. The extended t-isomorphism v preserves dagger exactly at every level. Applying lem-compcb-amplified-compression-identities to the square pairs (R\_n,R\_n) in A and (a\_n,a\_n) in A\_R shows successively that both compressions preserve dagger. Hence T\_n(x\textasciicircum{}dagger)=T\_n(x)\textasciicircum{}dagger for every n and x. \steptag{validated} \\[4pt]
\step{1.3} Uniform near-isometry. By the exact formula in node 1.2 and the forward estimate of lem-maincb-nested-corner-comparison specialized as in node 1.1.1, ||T\_n(x)-v\_n(x)||=||W\_n(v\_n(x))-v\_n(x)||<=C\_nest*t*||v\_n(x)||. Def-extended-delta-inclusion gives ||v\_n(x)||<=(1+t)||x||, so for t<=1 this is at most 2*C\_nest*t*||x||. Combining this with (1-t)||x||<=||v\_n(x)||<=(1+t)||x|| by the triangle and reverse-triangle inequalities gives (1-K\_N*t)||x||<=||T\_n(x)||<=(1+K\_N*t)||x|| with K\_N=1+2*C\_nest, uniformly in n. \steptag{validated} \\[4pt]
\step{1.4} Uniform homomorphism clauses. There are universal K\_M,K\_U such that, at every amplification, T\_n has multiplicative defect at most K\_M*t and unit defect at most K\_U*t; together with node 1.2 it is dagger preserving and linear. \steptag{validated} \\[6pt]
\step{1.4.1} Multiplicative telescope, conditional on nodes 1.1 and 1.3. Fix n and x,y in M\_n tensor B; put b=v\_n(x), c=v\_n(y), beta=T\_n(x), gamma=T\_n(y), and distinguish the original product of A by juxtaposition from the compressed R-product dot\_R, which is the ambient product of A\_R. The extended t-isomorphism clause gives ||v\_n(xy)-b dot\_P c||<=t||x||||y||. Lem-compcb-rectangular-product in A gives ||b dot\_P c-bc||<=2*C\_co*t||b||||c||. The ambient product bound and ||b||,||c||<=(1+t) times the source norms give a universal bound ||v\_n(xy)||<=K\_B||x||||y||, hence node 1.1.1 gives ||T\_n(xy)-v\_n(xy)||<=C\_nest*K\_B*t||x||||y||. Bilinearity of the original A-product gives bc-beta gamma=(b-beta)c+beta(c-gamma); the ambient product bound, node 1.3, and the uniform norm bounds control this by K\_pert*t||x||||y|| without associativity. Since beta,gamma lie in M\_n tensor S\_a\textasciicircum{}\{A\_R\} and A\_R=S\_R\textasciicircum{}A, lem-compcb-amplified-compression identifies them with elements of the amplified R,R-corner in A. A separate application of lem-compcb-rectangular-product in A therefore gives ||beta gamma-beta dot\_R gamma||<=2*C\_co*t||beta||||gamma||. In A\_R, whose ambient product is dot\_R and whose defect is at most 2*C\_ca*t, a is a C\_nest*t-projection by node 1.1.1; hence lem-compcb-rectangular-product applied inside A\_R gives ||beta dot\_R gamma-beta dot\_a gamma||<=C\_co*(C\_nest+2*C\_ca)*t||beta||||gamma||. Lem-compcb-amplified-compression identifies all amplified compressed products used here. Summing the six successive differences T\_n(xy)-v\_n(xy), v\_n(xy)-b dot\_P c, b dot\_P c-bc, bc-beta gamma, beta gamma-beta dot\_R gamma, and beta dot\_R gamma-beta dot\_a gamma gives ||T\_n(xy)-T\_n(x) dot\_a T\_n(y)||<=K\_M*t||x||||y|| for a universal K\_M independent of n and all dimensions. \steptag{validated} \\[6pt]
\step{1.4.1.1} Missing product bridge. At amplification n, beta,gamma are in M\_n tensor S\_a\textasciicircum{}\{A\_R\} by node 1.2. Because A\_R=S\_R\textasciicircum{}A, lem-compcb-amplified-compression identifies this space as a subspace of the amplified R,R-corner S\textasciicircum{}A\_\{R\_n,R\_n\}, where R\_n=I\_n tensor R. Apply lem-compcb-rectangular-product in A to the compatible R\_n,R\_n factors: since R is a t-projection and A is an extended t-C*-algebra, its total error is 2*t, so ||beta gamma-beta dot\_R gamma||<=2*C\_co*t||beta||||gamma||. Here dot\_R=Co\textasciicircum{}A\_\{R\_n,R\_n\}(beta gamma) is also exactly the ambient multiplication of M\_n tensor A\_R: the corner-algebra structure uses the compressed R-product, and tensor linearity together with lem-compcb-amplified-compression identifies its amplification with Co\textasciicircum{}A\_\{R\_n,R\_n\}. Separately, inside A\_R, node 1.1 gives algebra defect 2*C\_ca*t and projection defect C\_nest*t for a; the same external then yields ||beta dot\_R gamma-beta dot\_a gamma||<=C\_co*(C\_nest+2*C\_ca)*t||beta||||gamma||. Thus ||beta gamma-beta dot\_a gamma|| is bounded by C\_co*(2+C\_nest+2*C\_ca)*t||beta||||gamma||, supplying exactly the omitted intermediate term. \steptag{validated} \\[4pt]
\step{1.4.2} Unit telescope, conditional on nodes 1.1 and 1.3. Let u\_P=Co\textasciicircum{}A\_P(P) and u\_a=Co\textasciicircum{}\{A\_R\}\_a(a) be the compressed units. The t-homomorphism unit clause gives ||v\_n(1)-I\_n tensor u\_P||<=t. On M\_n tensor S\textasciicircum{}A\_P, node 1.1.1 gives ||W\_n||<=1+C\_nest*t, so the contribution of this unit error to ||T\_n(1)-I\_n tensor u\_a|| is at most (1+C\_nest*t)t. It remains to compare W\_n(I\_n tensor u\_P) with I\_n tensor u\_a. The t-projection estimate ||P||<=2, the compression-to-raw-map bound in def-compressed-corner, and ||P\textasciicircum{}2-P||<=t give ||u\_P-P||<=K\_0*t: for example use the raw association (P\textasciicircum{}2)P and the exact bilinear identity (P\textasciicircum{}2)P-P=(P\textasciicircum{}2-P)P+(P\textasciicircum{}2-P). The same compression comparison bounds Co\_R and Co\textasciicircum{}\{A\_R\}\_a uniformly because ||R||,||a||<=2. Since a=Co\_R(P) and u\_a=Co\textasciicircum{}\{A\_R\}\_a(a), exact linearity gives W\_1(u\_P)-u\_a=Co\textasciicircum{}\{A\_R\}\_a Co\_R(u\_P-P), and lem-compcb-amplified-compression tensors this identity to every n. Hence ||W\_n(I\_n tensor u\_P)-I\_n tensor u\_a||<=K\_1*t. Thus the unit defect is at most K\_U*t for universal K\_U. Together with node 1.2, the two children establish every homomorphism clause claimed by node 1.4. \steptag{validated} \\[4pt]
\step{1.5} Bijectivity. Node 1.1.1 makes W\_1 a linear map S\textasciicircum{}A\_P->S\textasciicircum{}\{A\_R\}\_a with ||W\_1(X)-X||<=C\_nest*t||X||. If C\_nest*t<1 and W\_1(X)=0, then ||X||<=C\_nest*t||X||, hence X=0, so W\_1 is injective. Lem-maincb-nested-corner-dimension-transport, specialized to Q=P as in node 1.1.1, gives equality of the finite dimensions of its domain and codomain; therefore W\_1 is surjective and hence bijective. Since v:B->S\textasciicircum{}A\_P is bijective by the assumed extended t-isomorphism, T=W\_1 composed with v is bijective. \steptag{validated} \\[4pt]
\step{1.6} Final assembly. After nodes 1.1--1.5 are validated, choose one positive universal e\_out below all comparison, compression, corner-algebra, rectangular-product, nonvanishing, and Neumann-free dimension thresholds used there, and choose C\_out=max\{1,C\_ca*(C\_nest+2*C\_ca),K\_N,K\_M,K\_U\}. Then every T\_n is a C\_out*t-inclusion, T is bijective, and the exact formula T\_n=I\_n tensor T holds, proving node 1. \steptag{validated} \\[4pt]
\end{proofbox}

\end{document}
